\section{Post-saturation revision-responsibility validation}
\label{sec:p1x}

The literature refresh after the v2.2.4 result raised a broader but distinct systems question: after diagnosis, repair, belief/model revision, dependency maintenance, and generic mutation authorization are all granted to strong donor mechanisms, does an explicit rule for \emph{which scientific layer is justified to change} add value? We tested that question in a separate prospectively frozen successor study, P1-X. This successor is not pooled with the v2.2.4 mechanical experiment and does not change its protocol, worlds, parents, or statistics.

\paragraph{Contract and comparators.}
P1-X separates three labels that need not coincide: causal responsibility (what produced the failure), best intervention (what would restore the task), and scientifically justified revision (what epistemic layer the evidence licenses changing). The candidate control contract requires registered strictly narrower admissible revisions to be insufficient before broader epistemic mutation, preserves protected invariants, binds the execution certificate to the observed reopen scope, and returns \texttt{REQUEST\_DISCRIMINATOR}, \texttt{UNRESOLVED}, or \texttt{CANNOT\_CHECK} rather than forcing an unsupported escalation. The comparator-fair V2 battery contains 400 exact cases across software/debugging, scientific retrieval, semantic integration, evidence/verification, and model/experiment-design families. Eight archetypes include narrow-repair sufficiency, high-level necessity, incomparable minima, diagnosis--prescription separation, external/evaluator responsibility, preservation failure, reopen-scope mismatch, and clean/insufficient-evidence controls. All arms receive matched candidate-visible information and intervention budgets.

The primary comparator B1 is a donor-complete greedy product: it receives the same diagnosis outputs, admissible candidates, protected checks, dependency information, generic authorization, and execution enforcement, but lacks the P1-X minimal scientific-escalation rule. B2 permits a successful high-level repair to authorize reframing after generic safety checks. B3 is deliberately stronger: an ideal information-equivalent product receives the exact P1-X responsibility classes, partial order, counterfactual results, preservation/reopen predicates, and authority semantics. If those semantics are identical, B3 should tie P1-X; it is an expressivity boundary rather than a baseline that P1-X is expected to defeat.

\paragraph{Negative V1 retained and disjoint V2 repair.}
The first protected execution passed its literal frozen statistical gates, but an independent post-access reconstruction found a comparator implementation defect: the B1/B2 clean \texttt{NO\_CHANGE} guard iterated dictionary keys instead of status values. We did not patch that experiment after observing outcomes. V1 remains archived as positive-but-comparator-compromised evidence. A disjoint V2 then prospectively fixed only that comparator bug, preserved the scientific hypothesis, P1-X/B3 semantics, primary $+0.10$ practical margin and non-regression rules, and used new protected identities with seed-bound proposal ordering.

\paragraph{Comparator-fair V2 result.}
On the 400 V2 cases, P1-X obtains 400/400 exact scientific revision decisions (ESRD $=1.000$), compared with 275/400 ($0.6875$) for B1 and 200/400 ($0.500$) for B2. The paired P1-X--B1 difference is $+0.3125$ with a predeclared domain-stratified bootstrap 95\% interval $[0.2675,0.3575]$, exceeding the $+0.10$ practical margin; exact McNemar counts are $b=125,c=0$. P1-X records zero false high-level reframes and zero protected-invariant violations, and its advantage over B1 is positive in each of the five domains. An independent implementation reproduces all 400 identities, the protected-bundle and decision-row hashes, every headline arm count, and zero B3 decision mismatches.

B3 obtains 400/400 and is decision-identical to P1-X. This negative boundary is central: the result does \emph{not} show an inherent expressivity or centralization advantage. It shows that making the scientific escalation contract explicit can matter relative to a realistic donor-complete interface that has the component mechanisms but not the same coupling; an ideal product supplied identical semantics is equivalent.

\paragraph{Interface-substitution check.}
A separate exact 200-case phase tested four equally informative diagnosis/recovery ordering combinations. P1-X remains 200/200 in all four combinations, B3 remains exactly identical, and P1-X exceeds B1 in every combination and every domain with zero false reframes, invariant violations, or budget violations. This supports bounded interface-order invariance. It is not evidence that P1-X has been validated over independently trained diagnosis/recovery engines or deployed open-ended scientific agents.

The authorized successor claim is therefore limited to the exact heterogeneous contract families and interface substitutions tested here. Open-ended/model-bearing scientific-agent generality remains \texttt{CANNOT\_CHECK}.
